\documentclass[cn,hazy,blue,12pt,normal]{mynote}
\begin{document}

\title{物理化学教学研讨会}
\author{张斌}
\maketitle

\section{杨学明\ 物理化学中的反应动力学}

物理与化学的交叉是推动化学学科发展的重要推动力。

\begin{figure}[htpb]
    \centering
    \includegraphics[width=.6\linewidth]{杨学明}
\end{figure}

介绍一本物理化学教材 Physical Chemistry - A Molecular Approach

\begin{itemize}
  \item 第一章\ 量子理论的黎明
  \item 第二章\ 经典波动方程
  \item 第三章\ 薛定谔方程
  \item $\cdots$
  \item 二十八章\ 宏观反应动力学\ 速率定律
  \item 二十九章\ 宏观反应动力学\ 反应机理
  \item 三十章\ 微观反应动力学
  \item 三十一章\ 固体与表面化学
\end{itemize}

现代物理化学的基本逻辑：

\begin{itemize}
  \item 微观结构机理$\rightarrow$宏观化学特性
  \item 化学科学根本是分子的结构及化学变化
\end{itemize}

过去二十多年中，化学反应动力学研究得到了快速发展！

化学动力学研究对于化学科学发展具有重要意义！

微观反应动力学范例：

\begin{itemize}
  \item 基元化学反应机理
    \begin{itemize}
        \item 多通道化学反应动力学
        \item 插入醒化学反应机理
        \item 萃取型化学反应及其量子效应
        \item 受困萃取型化学反应
        \item 共振型的化学反应 \ch{F + H2 -> HF + H}
    \end{itemize}
  \item 化学反应共振
  \item 化学反应中的几何相位效应
  \item 分子光化学反应动力学
\end{itemize}

化学动力学研究的里程碑：

\begin{itemize}
  \item 阿仑纽斯方程：1903年Nobel化学奖
  \item 过渡态理论：1980年Nobel化学奖
  \item 分子反应动力学研究的突破：1986年Nobel化学奖
  \item 电子转移及RRKM理论的发展：1992年Nobel化学奖
  \item 臭氧层破坏机理的发现：1995年Nobel化学奖
  \item 分子束中C60分子的发现：1996年Nobel化学奖
  \item 量子化学理论的发展：1998年Nobel化学奖
  \item 超快反应动力学的突破：1999年Nobel化学奖
  \item 表面催化反应机理的发展：2007年Nobel化学奖
  \item 大分子动力学理论的发展：2013年Nobel化学奖
\end{itemize}

高灵敏度的实验方法以及精确的动力学理论是原动力！

化学动力学理论研究进入精确量子动力学时代。

薛定谔方程是化学量子动力学理论的基础，针对不同反应体系的薛定谔方程的精确求解是化学动力学理论发展的核心问题。

对于实际化学体系的重要影响：
\begin{itemize}
  \item 大气化学
  \item 燃烧化学
  \item 化学激光
  \item 星际化学
\end{itemize}







\newpage

\section{陈军\ 超越传统的电池体系}

\begin{figure}[htpb]
    \centering
    \includegraphics[width=.6\linewidth]{陈军}
\end{figure}

立项2个亿

储能在新型电力系统中处于关键地位

\begin{itemize}
  \item 物理储能
  \item 化学储能---电池
\end{itemize}

二次电池能可逆充放电，在储能领域具有重要作用。

储能电池需求量快速飙升。

电池科技的发展将有力推动汽车的电动化进程。

\begin{figure}[htpb]
    \centering
    \includegraphics[width=\linewidth]{效率}
\end{figure}

汽车电动化的必要性：
\begin{itemize}
  \item 纯电动车：77\%
  \item 氢燃料电池车：30\%
  \item 合成染料内燃机汽车：13\%
\end{itemize}



超越传统的电池体系：
\begin{itemize}
  \item 1859 法国\ 铅酸电池 35-50 Wh/kg
  \item 1970-1990s 美国，日本\ 锂离子电池 150-250 Wh/kg
  \item 未来电池：新理论，新机制，新材料，新范式
\end{itemize}

四大现存挑战：
\begin{itemize}
  \item 安全
  \item 电化学性能逼近理论极限
\end{itemize}



\newpage

\section{江雷\ 董智超\ 对基础课物理化学中表界面化学教学内容改革的一些意见}

\begin{figure}[htpb]
    \centering
    \includegraphics[width=.6\linewidth]{董智超}
\end{figure}

界面几乎是所有功能的源泉。

界面与浸润性：

\begin{itemize}
  \item 思考一：与日常相关较强
  \item 思考二：比表面自由能的数值可以用简单的方法算出来吗？
  \item 思考三：由于表面张力是一个线张力，这个定义是否可以直接变为宽度，方便今后学习和使用中的理解。
\end{itemize}

界面超浸润现象：

\begin{itemize}
  \item 荷叶表面超疏水
  \item 动物角膜具有超亲水性
\end{itemize}

杨氏方程：所有液体亲属界限为90°？不一定

光学测量接触角的方法存在局限，在超疏水表面有很大误差。

\newpage

\section{张玮\ 红色基因，踔厉奋发------北京理工大学课程思政津设索与实践}

\begin{figure}[htpb]
    \centering
    \includegraphics[width=.6\linewidth]{张玮}
\end{figure}

党的二十大报告对教育的要求

党和国家对课程思政的要求











\newpage
\section{原弘\ 精讲“相平衡”}

\begin{figure}[htpb]
    \centering
    \includegraphics[width=.6\linewidth]{原弘}
\end{figure}

相平衡知识背景

相平衡基本内容：

\begin{itemize}
  \item 相律/相图
  \item 
\end{itemize}

相平衡讲解方法

黄子卿：水的三相点












\newpage

\section{白守礼\ 精讲“热力学第二定律”}

\begin{figure}[htpb]
    \centering
    \includegraphics[width=.6\linewidth]{白守礼}
\end{figure}

中国人历史上对热力学第二定律的描述：

\begin{itemize}
  \item 水往低处流
  \item 黄河之水天上来，奔流到海不复回。
\end{itemize}












\newpage
\section{刘向荣\ 精讲“化学动力学”}

教学思考

\begin{itemize}
  \item 课堂教学的作用是什么
  \item 课堂听讲的作用是什么
\end{itemize}

互动式教学策略

\begin{itemize}
  \item 教师先动：课前设计问题，发布在qq群里，鼓励同学们在预习中考虑并给出自己的看法
  \item 学生后动：新开课的班，先给大家示范，如何在书上找问题答案，让学生们动起来
  \item 师生互动：课上抛出课前提的问题，鼓励同学们积极响应，给出自己的看法，教师和同学们共同讨论判断正误
\end{itemize}

抛出问题：

\begin{itemize}
  \item 热力学和动力学的异同？
  \item 同学们能用时间轴画出化学动力学的发展历程吗？
  \item 为何在化学反应速率的表达式中引入转化速率($\mathrm{d}\xi /\mathrm{d}t$)？
  \item 一级反应速率方程建立的方法是什么？
  \item 具有简单级数反应速率方程的特征是什么？
\end{itemize}



\newpage
\section{朱志昂\ 我国物理化学教材建设回顾和40年教材建设的经验体会}

教学不要停留在课本

如何用现在的观念讲以前的东西

放眼世界，站得更高

\begin{figure}[htpb]
    \centering
    \includegraphics[width=.6\linewidth]{朱志昂}
\end{figure}

一、我国物理化学教材建设历史的回顾：

1. 1953年暑期，教育部在青岛召开综合性大学化学系
教学计划讨论会，参照苏联五年制教学计划制定了我
国自己的教学计划，把原先的物理化学课程一分为四
，即物理化学（约110学时）、物质结构（约50学时）
、胶体化学（约30学时）和结晶化学（约30学时）。

2，1954年2月，在北京大学开会讨论黄子卿先生拟定的
教材教学大纲，这是我国自行拟定的第一个物理化学
教学大纲。黄子卿先生根据这个大纲，于1955年出版
了《物理化学》一书，这是我国自行编写出版的第一
部物理化学教材


3. 1959年，华东化工学院院长张江树教授主编了《物
理化学与胶体化学》。

4. 1961年6月教育部在南京召开化学系几门基础课教
的选编会议，会上决定以南京大学的讲义为基础，适
当补充一些章节，这就是南京大学傅献彩先生的《物
理化学》第一版。

二、教材编写中应该注意的几个问题
\begin{itemize}
  \item 教材与授课内容的关系
  
  教材是教学的基本依据。如果没有教材，学生在课后
  就没有主要的参考书，无所依从。另一方面，讲课没
  有必要与教材完全吻合，一切按教材讲授，教材上写
  多少就讲多少，就失去了自己的特点。
  \item 培养创新性人才，首先要培养学生的质疑精神和自
  学能力。
  
  学生有疑问可去看书得到解决，有兴趣的学
  生还可看老师课堂未讲的拓展性的内容，教材一定要
  比课堂讲授的内容丰富一些，课堂上要讲重点难点，
  还要讲清楚，但要留一部分让学生自学，这些部分要
  通过考试或习题来考核。

  \item 在教材和讲课之间还普通存在一种现象，就是根据学
  时数来选择讲授教材中的部分章节，学时少了，有的
  整章就不讲。这一作法是不科学的。
  
  早在20世纪八十
  年代初，吴征铠教授就指出：学生在高等学校学习不
  过四、五年，而在工作中的学习却是四、五十年，如
  何在四、五年中为以后的四、五十年做好准备，是师
  生必须共同牢记的，教师更应反复说明的问题。

  \item 我们应该让学生掌
  握物理化学各章的基本原理和基本方法，将来在工作
  中遇到问题时有一基本的思路。
  
  在这一思想指导下，
  2017年、2018年秋季，西安交通大学钱学森学院聘我
  为化学生物学专业本科生讲授物理化学课程，选用的
  教材是南京大学傅献彩的《物理化学》讲授64学时，
  我采用各章都讲的方法。

  \item 教材建设和科研教研的关系:

  教材内容是生产实践和科学实验规律的总结和提炼。
  物理化学各章均为各三级学科的基本研究方法，是之
  前科研的总结提升，也是学生未来开展科研的基础。
  根据我的亲身经历，在教学和科研取得一定成果后再
  编写教材，可以得心应手，水到渠成。

  \item 科研和教研的加持使教材的编写更加鲜活，水平也不
  断提升。
  
  物理基本原理在科研中的应用，
  将教学研究成果直接融入教材。例如“化学平衡教学
  中的某些问题”、“热力学标准态和化学反应的标准
  热力学函数”、“关于电池反应方向判据的讨论”
  “关于电动电势计算公式的讨论"、“讲授物理化学
  律的探讨”
\end{itemize}

三、对未来教材建设的建议

\begin{enumerate}
  \item 注重宏观和微观的结合
  \item 压缩经典，不断增加学科前沿知识
  \item 与课程思政教育有机融合
  \item 教材建设的多元化、立体化
    \begin{itemize}
      \item 前沿拓展
      \item 建设信息化资源
    \end{itemize}
\end{enumerate}


思考：

如何从现代科学技术的视觉学习和讲授一百多年前
物理化学的定律定理？如何融入现代科学技术成果


\newpage
\section{彭笑刚\ 大学化学与物理化学}

\begin{figure}[htpb]
    \centering
    \includegraphics[width=.6\linewidth]{彭笑刚}
\end{figure}

《大学化学讲义》 彭笑刚

\subsection{基本事实与基本判断}

\begin{itemize}
  \item 过去一百年，化学积累了绝大部分方法、知识和智慧
    \begin{itemize}
        \item 确定元素的实在性，以及天然元素的全部数目
        \item 确定基本粒子（原子）的基本结构与性质：波粒二象性的电子+核
        \item 确定原子内与原子间相互作用：量子化的电子相互作用
        \item 系统且定量的化学全域框架大体成型，初步完成波义尔使命
        \item 化学第二使命动能充足，在原子---分子水平重构物质成为常态
    \end{itemize}
  \item 脱胎于苏联系统的教学体系不满足养育中国年轻一代的需求
  \item 年轻学子精神与心理健康状态堪忧
    \begin{itemize}
      \item “北大40\%的大学生应该看心理医师”
      \item 浙大物理化学课堂至少30\%的学生需要心理辅导
      \item 重塑国民性不是一个伪命题，而是迫切需求
      \item 农耕文明教育观与学习观不能让它再继续下去了
      \item 帮助学生确立“大学学习是寻找实践人生方式的优良途径”
      \item 科学（化学）是开放系统、有生命的系统，可以成为实践人生的精彩途径
    \end{itemize}
\end{itemize}

\subsection{通识课、基础课与专业课}

\begin{itemize}
  \item 一门通识课：大学化学(或普通化学)
  \item 三门基础课：有机化学、结构化学、物理化学
  \item 几门专业课：中级无机、中级分析化学、高分子化学、化学生物学
  \item 三门实验课：基础化学实验、合成化学实验、化学测量学实验
  \item 高级课程：本研贯通课程
\end{itemize}

\subsection{三个重要概念}

\begin{itemize}
  \item 熵 $S = k \ln W$
  \item 温度 $T \equiv  (\frac{\partial V}{\partial T})_{V,N}$
  \item 内能
\end{itemize}


\newpage

\section{侯文华\ 国外物理化学教学内容更新的现状分析}

\begin{figure}[htpb]
    \centering
    \includegraphics[width=.6\linewidth]{侯文华}
\end{figure}

三本书：

\begin{itemize}
  \item Atkins物理化学第11版
  \item Levine 编写的 Physical Chemistry (6th edition) 
  \item Physical Chemistry - A Molecular Approach
\end{itemize}


\newpage

\section{张韫宏\ 雾霾形成机制的几个热力学和动力学问题}

气溶胶与大气颗粒物

大气圈三次危机：

\begin{itemize}
  \item 臭氧洞
  \item 温室效应
  \item 雾霾
\end{itemize}

二次气溶胶形成的复杂物理化学过程

颗粒物的酸碱平衡问题

抛出三个问题：

\begin{enumerate}
  \item 教师培养的问题
  \item 教学改革的问题
  \item AI的问题
\end{enumerate}


\newpage

\section{郭东升\ 物理化学基础识在科学前沿中的应用}

\begin{figure}[htpb]
    \centering
    \includegraphics[width=.6\linewidth]{郭东升}
\end{figure}

研究方向：生物医用超分子材料

\begin{itemize}
  \item 基础：分子设计合成
  \item 核心：识别组装性质
  \item 导向：生物医药功能
\end{itemize}

高水平教学促进高水平科研；高水平科研反哺高水平教学。


\newpage

\section{陈亚芍\ 课堂讲授与慕课混合模式的教学探讨}

\begin{figure}[htpb]
    \centering
    \includegraphics[width=.6\linewidth]{陈亚芍}
\end{figure}

\begin{itemize}
  \item 混合式教学的内涵与现状
  \item 混合式教学的有效策略思考
    \begin{itemize}
      \item 混合式教学的关键是对课程进行再设计
      \item 已有MOOC课程资源与课堂教学结合组织混合式教学
      \item 利用SPOC教学平台组织混合式教学
      \item 提高教师的信息化技术与教育教学深度融合能力，推动混合式教学发展
      \item 形成性评价推动混合式教学发展
    \end{itemize}
  \item 形成性评价探索
  \item 混合式教学的实践探索
  \item 总结与展望
\end{itemize}


black board；智慧树平台



\end{document}